% Generated draft: requires mathematical and source review.
\documentclass[11pt]{article}
\usepackage[margin=25mm]{geometry}
\usepackage{fontspec}
\setmainfont{Latin Modern Roman}
\usepackage{amsmath,amssymb,mathtools,unicode-math}
\setmathfont{Latin Modern Math}
\usepackage{enumitem,needspace,xurl,hyperref}
\setlength{\parindent}{0pt}
\setlength{\parskip}{5pt}
\setlength{\emergencystretch}{3em}
\newcommand{\N}{\mathbb{N}}
\newcommand{\Z}{\mathbb{Z}}
\newcommand{\Q}{\mathbb{Q}}
\newcommand{\R}{\mathbb{R}}
\newcommand{\C}{\mathbb{C}}
\newcommand{\F}{\mathbb{F}}
\newcommand{\A}{\mathbb{A}}
\newcommand{\PP}{\mathbb{P}}
\newcommand{\E}{\mathbb{E}}
\newcommand{\eps}{\varepsilon}
\newcommand{\dd}{\,\mathrm{d}}
\begin{document}
\section*{UnsolvedMath ID: 30004056}
\textbf{Source problem code:} catalogue entry
\subsection{Definitions and mathematical statement}
Let \((G,X)\) denote a Shimura datum, where \(G\) is a connected reductive group over \(\Q\), and let \(E=E(G,X)\) be its reflex field. For a compact open subgroup \(K\subset G(\A_f)\), write \(\operatorname{Sh}_K(G,X)\) for the Shimura variety over \(E\); after choosing an embedding \(E\hookrightarrow\C\), its complex points have the description
\[
\operatorname{Sh}_K(G,X)(\C)
=G(\Q)\backslash\bigl(X\times G(\A_f)/K\bigr),
\]
where \(\A_f\) denotes the finite adeles of \(\Q\). Parahoric level at a prime \(p\) refers to a factorization \(K=K_pK^p\), where \(K_p\subset G(\Q_p)\) is a parahoric subgroup in the sense of Bruhat--Tits theory and \(K^p\subset G(\A_f^p)\) is a compact open subgroup away from \(p\). The precise convention for the parahoric subgroup is not specified in the supplied record.

To display the missing formulation choices, suppose a place \(v\mid p\) of \(E\) is chosen and put \(B=\mathcal O_{E,v}\). In a scheme formulation, a flat integral model over \(B\) consists of a flat \(B\)-scheme \(\mathscr S\) together with an identification
\[
\iota:\mathscr S\times_B E\xrightarrow{\sim}\operatorname{Sh}_K(G,X).
\]
This describes one possible local framework; the supplied record does not choose this base, the category of models, or additional conditions such as finite type, separatedness, normality, or an extension property.

\textbf{Statement requested by the source.} Determine the existence of integral models of Shimura varieties with parahoric level structure, and characterize their uniqueness. A precise version would have to specify the allowed Shimura data, primes, levels, bases, and a defining property \(\mathcal P\) of the models, then determine for which allowed data a pair \((\mathscr S,\iota)\) satisfying \(\mathcal P\) exists and in what sense it is unique. For example, uniqueness up to an isomorphism compatible with the generic-fibre identifications asks, for two such models, for a \(B\)-isomorphism \(f:\mathscr S_1\to\mathscr S_2\) with
\[
\iota_2\circ f_E=\iota_1;
\]
uniqueness of this compatible isomorphism is a further possible requirement. The source does not distinguish these meanings or supply \(\mathcal P\). Thus the supplied question is a broad existence-and-characterization target, not a fully quantified assertion of existence or uniqueness for every Shimura datum.
\subsection{Short English statement}
The problem asks when Shimura varieties with parahoric level admit integral models and which conditions characterize those models uniquely. The supplied statement leaves the permitted data, base, category of models, defining properties, and meaning of uniqueness unspecified, so it requires clarification before it can become a single precise conjecture.
\subsection{Sources}
\begin{itemize}
\item Supplied local UnsolvedMath registry record, permanent ID \(30004056\), title \emph{Integral Models of Parahoric Shimura Varieties}, source problem code \texttt{catalogue entry}, in the Oberwolfach Reports Open Problems collection. No individual proposer is named; the supplied proposed year is 2019.
\item Supplied report citation: \emph{Arithmetic of Shimura Varieties} (2020), Oberwolfach Reports, pp.~65--131, DOI \texttt{10.4171/owr/2019/2}. \url{https://doi.org/10.4171/owr/2019/2}. The background attributes the record to workshop open-problem discussions documented in Oberwolfach Reports; no specific page within this range is supplied.
\item Kisin--Pappas, ``parahoric integral models'' (supplied reference label). \url{https://arxiv.org/abs/1512.01149}.
\item Kisin--Pappas--Zhou, ``integral models II'' (supplied reference label). \url{https://arxiv.org/abs/2409.03689}.
\item Literature assessment dated exactly \texttt{2026-08-22}: the supplied curation labels the problem open and describes it as a broad existence, characterization, and uniqueness question. It reports constructions in major cases and states that no general uniqueness-and-characterization result resolves the stated problem. This is supplied dated curation, not an independently verified literature assessment; its time-sensitive status and cited claims require independent verification.
\item The registry supplies the statement-status label \texttt{corrected\_verified} and reports a manual semantic, typographic, and LaTeX review. These are supplied provenance claims, not independent confirmation that the mathematical target is fully specified.
\end{itemize}
\end{document}
